\documentclass[11pt]{article}
\usepackage[margin=0.9in]{geometry}
\usepackage{amsmath,amssymb,amsthm}
\usepackage[hidelinks]{hyperref}
\newtheorem{theorem}{Theorem}[section]
\newtheorem{lemma}[theorem]{Lemma}
\newtheorem{proposition}[theorem]{Proposition}
\newtheorem{corollary}[theorem]{Corollary}
\newcommand{\F}{\mathbb F}
\newcommand{\PP}{\mathcal P}
\title{Critical-scale growth bounds for congruence-preserving functions over finite fields}
\author{Henry Zweiman}
\date{September 9, 2026}
\begin{document}
\maketitle
\begin{abstract}
Let $R=\F_q[t]$. We prove that a function $f:R\to R$ preserving congruences
modulo every irreducible polynomial is a polynomial map if
$\deg f(A)\le q^{\deg A}/(256q^2)$ for every input of sufficiently large
degree. This replaces the factor $q^n/n$ in the growth hypothesis of
Bell and Nguyen by a positive constant times $q^n$. The proof uses
integer-valued Carlitz digit polynomials, multiplied by a squarefree
primorial, in the auxiliary-polynomial method. The rationality step is
Bell and Nguyen's linear-growth theorem. For $\F_q$-linear functions we
prove the sharp threshold $q/(q-1)$ for
$\limsup_n\deg f(t^n)/q^n$, and give a nonpolynomial example attaining it.
We also give a Newton-coefficient formulation of the remaining sharp
problem and an obstruction to a coefficient estimate using growth alone.
\end{abstract}

\section{Results and their scope}
Throughout, $q$ is a prime power, $R=\F_q[t]$, and $K=\F_q(t)$.
We set $\deg 0=-\infty$. Write $V_n=\{A\in R:\deg A<n\}$, including
zero, and let $\PP$ be the set of monic irreducible polynomials in $R$.
We say that $f:R\to R$ preserves prime congruences if
\begin{equation}\label{eq:cp}
 P\mid A-B\quad\Longrightarrow\quad P\mid f(A)-f(B)
 \qquad(A,B\in R,\ P\in\PP).
\end{equation}
A polynomial map is the evaluation of one polynomial in $K[X]$ on $R$.
Its coefficients need not belong to $R$.

Bell and Nguyen \cite[Theorem 1.4]{BN} proved polynomiality under
$\deg f(A)<q^n/(27qn)$ for all sufficiently large $n=\deg A$.
Their Section 4 asks whether the growth function $q^n/n$ can be enlarged;
Question 4.1 proposes every fixed constant below the primorial threshold.
Our first theorem addresses the enlargement of the growth scale.

\begin{theorem}\label{thm:general}
Let $\ell$ be the least positive integer such that $q^\ell\ge16$, and set
\[
 c_q=\frac1{16q^{\ell+1}}.
\]
If $f:R\to R$ satisfies \eqref{eq:cp} and
\begin{equation}\label{eq:growth}
 \deg f(A)\le c_q q^{\deg A}
\end{equation}
for every input of sufficiently large degree, then $f$ is a polynomial map.
In particular, the bound $q^{\deg A}/(256q^2)$ suffices.
\end{theorem}

\begin{corollary}\label{cor:littleo}
If $f$ satisfies \eqref{eq:cp} and
$\max_{\deg A=n}\deg f(A)=o(q^n)$, then $f$ is a polynomial map.
\end{corollary}

The constant in Theorem \ref{thm:general} is chosen for a transparent
proof, rather than optimized. For $q=2,3,4,9$ the stated $c_q$ is,
respectively, $1/512,1/1296,1/1024,1/11664$.
Put
\begin{equation}\label{eq:primorial}
 \Pi_n=\prod_{\substack{P\in\PP\\\deg P\le n}}P,
 \qquad d_n=\deg\Pi_n,\qquad \Pi_0=1.
\end{equation}
Then $d_n/q^n\to q/(q-1)$. The proposed sharp condition in
\cite[Question 4.1]{BN} is $\deg f(A)\le(1-\varepsilon)d_n$ on every
sufficiently large degree-$n$ shell, for some $\varepsilon>0$.
Theorem \ref{thm:general} does not prove that full assertion.

\begin{theorem}\label{thm:additive}
Suppose $f:R\to R$ is $\F_q$-linear and satisfies \eqref{eq:cp}.
If $f$ is not a polynomial map, then
\[
 \limsup_{n\to\infty}\frac{\deg f(t^n)}{q^n}\ge\frac q{q-1}.
\]
This constant is sharp, even with the maximum over all inputs of degree
$n$ in the numerator. Thus Bell--Nguyen Question 4.1 holds for
$\F_q$-linear functions.
\end{theorem}

The interpolation polynomials below are classical. Wagner's work
\cite{W74,W76} describes related coefficient divisibility for functions
preserving all polynomial moduli, including prime powers. Here only
\eqref{eq:cp} is assumed, and the squarefree factor $\Pi_n$ replaces the
least common multiple with prime powers. Li and Sha \cite{LS} study
congruence-preserving maps between finite residue rings. We use the
classical interpolation tools to obtain a growth theorem for arbitrary
maps on the infinite ring $R$.

\section{Integral interpolation polynomials}
Define
\begin{equation}\label{eq:carlitz}
 e_i(X)=\prod_{B\in V_i}(X-B),\qquad
 H_i=e_i(t^i),\qquad E_i(X)=\frac{e_i(X)}{H_i}.
\end{equation}
In particular $e_0(X)=E_0(X)=X$ and $H_0=1$.
These are the normalized linear Carlitz polynomials; see also
\cite[Section 2]{W74}. We give the needed elementary properties.

\begin{lemma}\label{lem:carlitz}
The polynomial $E_i$ is $\F_q$-linear, has $X$-degree $q^i$, and takes
values in $R$ on $R$. Moreover,
\begin{align}
 \deg H_i&=iq^i,\label{eq:Hdegree}\\
 v_P(H_i)&=\sum_{\substack{k\ge1\\k\deg P\le i}}q^{i-k\deg P},
       \label{eq:Hvaluation}\\
 E_i(t^j)&=0\ (j<i),\qquad E_i(t^i)=1,\label{eq:triangular}\\
 \deg E_i(A)&=(\deg A-i)q^i\quad(\deg A\ge i).
       \label{eq:Edegree}
\end{align}
\end{lemma}
\begin{proof}
The decomposition $V_{i+1}=V_i+\F_qt^i$ gives inductively
\[
 e_{i+1}(X)=e_i(X)^q-H_i^{q-1}e_i(X).
\]
This proves linearity and the asserted $X$-degree. The factors of $H_i$
are exactly all monic degree-$i$ polynomials, proving \eqref{eq:Hdegree}.
For $P$ of degree $d$, the number of those polynomials divisible by
$P^k$ is $q^{i-kd}$ if $kd\le i$ and zero otherwise. This proves
\eqref{eq:Hvaluation}.

For arbitrary $A$, if $e_i(A)=0$ there is nothing to prove. Otherwise,
for every $k$ with $kd\le i$, the map $V_i\to R/(P^k)$ is onto and has
fibers of size $q^{i-kd}$. Therefore
$v_P(e_i(A))\ge v_P(H_i)$. This holds at every prime, so $E_i(A)\in R$.
The triangular assertions follow from the definitions. If $\deg A\ge i$,
every factor $A-B$, $B\in V_i$, has degree $\deg A$, proving
\eqref{eq:Edegree}.
\end{proof}

\begin{lemma}\label{lem:primecount}
For $n\ge1$,
\[
 q^n\le d_n<\frac q{q-1}q^n,
 \qquad
 d_n=\frac{q(q^n-1)}{q-1}+O_q(q^{n/2}).
\]
\end{lemma}
\begin{proof}
The squarefree polynomial $t^{q^n}-t$ divides $\Pi_n$, proving the lower
bound. If $I_j$ counts monic irreducibles of degree $j$, then
$jI_j\le q^j$, whence $d_n\le\sum_{j=1}^nq^j<q^{n+1}/(q-1)$.
From $jI_j=\sum_{a\mid j}\mu(a)q^{j/a}$, the error in replacing $jI_j$
by $q^j$ is at most $\sum_{r\le j/2}q^r=O_q(q^{j/2})$.
Summing gives the final assertion.
\end{proof}

For $0\le j<q^s$, write $j=\sum_{i=0}^{s-1}j_iq^i$ with $0\le j_i<q$.
Set
\begin{equation}\label{eq:digit}
 F_j(X)=\prod_{i=0}^{s-1}E_i(X)^{j_i},\qquad
 G_{s,j}(X)=\Pi_{s-1}F_j(X).
\end{equation}
The first expression is independent of appended zero digits; $F_0=1$.

\begin{lemma}\label{lem:digit}
The polynomials $G_{s,j}$, $0\le j<q^s$, are $R$-valued and preserve
prime congruences. Their $X$-degrees are $j$, so they are linearly
independent over $K$. Put $\kappa=(q+1)/(q-1)$. For $D\ge s-1$,
\begin{equation}\label{eq:digitheight}
 \deg G_{s,j}(A)<(D-s+\kappa)q^s
 \qquad(A\in V_{D+1}).
\end{equation}
\end{lemma}
\begin{proof}
Integrality follows from Lemma \ref{lem:carlitz}. If $\deg P\le s-1$,
every value of $G_{s,j}$ is divisible by $P$. If $\deg P>s-1$, all the
$H_i$ in \eqref{eq:digit} are $P$-units. Thus $G_{s,j}\in R_{(P)}[X]$
and preserves congruences modulo $P$. Also $\deg_X F_j=j$.

If a factor with $i>\deg A$ has positive exponent, the value is zero.
In every other case, Lemmas \ref{lem:carlitz} and \ref{lem:primecount}
give, with $C=q/(q-1)$,
\begin{align*}
 \deg G_{s,j}(A)
 &\le d_{s-1}+\sum_{i=0}^{s-1}(q-1)(D-i)q^i\\
 &<\frac{q^s}{q-1}+(D-s+C)q^s-D-C
 <(D-s+\kappa)q^s.
\end{align*}
Here the same strict bound for $d_0=0$ applies when $s=1$.
\end{proof}

\section{The general growth theorem}
We first recall the elementary vanishing principle underlying the
Bell--Nguyen polynomial method \cite[Lemma 2.1]{BN}.

\begin{lemma}\label{lem:vanishing}
Suppose $g:R\to R$ satisfies \eqref{eq:cp}, vanishes on $V_{M+1}$, and
$\deg g(A)<d_{\deg A}$ whenever $\deg A>M$. Then $g=0$.
\end{lemma}
\begin{proof}
If not, choose $A$ of least degree $D$ with $g(A)\ne0$. Then $D>M$.
For each prime $P$ of degree at most $D$, the remainder of $A$ modulo
$P$ has degree less than $D$ and hence has zero image. Condition
\eqref{eq:cp} gives $\Pi_D\mid g(A)$, contradicting its degree bound.
\end{proof}

\begin{proposition}\label{prop:curve}
Under the hypotheses of Theorem \ref{thm:general}, there is a nonzero
$Q\in K[X,Y]$ such that $Q(A,f(A))=0$ for every $A\in R$.
\end{proposition}
\begin{proof}
Take $M$ sufficiently large and put
\[
 s=M-\ell\ge1,\qquad T=\lfloor q^M/4\rfloor,
 \qquad J=4q^{\ell+1}.
\]
We may require $\deg f(A)\le c_qq^M$ for every $A\in V_{M+1}$:
the finitely many exceptions to \eqref{eq:growth} have bounded image
degrees. Choose unknowns $u_{h,j,k}\in\F_q$ and form
\begin{equation}\label{eq:auxiliary}
 Q(X,Y)=\sum_{h=0}^T\sum_{j=0}^{q^s-1}\sum_{k=0}^J
       u_{h,j,k}t^hG_{s,j}(X)Y^k.
\end{equation}
Then $g(A)=Q(A,f(A))$ is $R$-valued and preserves prime congruences.
Since $\ell\le4$, $\kappa\le3$, $q^\ell\ge16$, and $Jc_q=1/4$,
Lemma \ref{lem:digit} gives
\begin{equation}\label{eq:ballheight}
 \deg g(A)<\frac{q^M}{4}+(\ell+\kappa)q^{M-\ell}
                       +\frac{q^M}{4}
 \le\frac{15}{16}q^M<q^M
 \quad(A\in V_{M+1}).
\end{equation}
Thus requiring $g(A)=0$ for these $q^{M+1}$ inputs imposes at most
$q^{2M+1}$ homogeneous linear equations over $\F_q$. The number of
unknowns is strictly larger:
\[
 (T+1)q^{M-\ell}(J+1)
 >\frac{q^M}{4}q^{M-\ell}4q^{\ell+1}=q^{2M+1}.
\]
A nonzero solution exists. It represents a nonzero polynomial $Q$:
for each power of $Y$, the $G_{s,j}$ have distinct $X$-degrees, and
$1,t,\ldots,t^T$ are independent over $\F_q$.

For an input of degree $D=M+r$, $r\ge0$, the same calculation gives
\begin{equation}\label{eq:outerheight}
 \frac{\deg g(A)}{q^D}
 <\frac{q^{-r}}4+(r+\ell+\kappa)q^{-\ell-r}+\frac14
 \le\frac{15}{16}.
\end{equation}
The two $r$-dependent summands are nonincreasing: for $a\ge1$,
$(r+a+1)/q\le r+a$. By Lemma \ref{lem:primecount} we have
$q^D\le d_D$. Lemma \ref{lem:vanishing} therefore gives $g=0$.
\end{proof}

\begin{proof}[Proof of Theorem \ref{thm:general}]
Clear denominators in the relation from Proposition \ref{prop:curve},
and write
\[
 Q=\sum_{k=0}^vP_k(X)Y^k,\qquad P_k\in R[X],\quad P_v\ne0.
\]
Necessarily $v>0$, since a nonzero polynomial in $X$ cannot
vanish on all of $R$. There are constants $a,b\ge0$ such that
\[
 \deg P_k(A)\le a\deg A+b
\]
for every nonzero $A$ and every $k$.
If $P_v(A)\ne0$ and $\deg f(A)>a\deg A+b$, the term
$P_v(A)f(A)^v$ has strictly larger degree than every other term of
$Q(A,f(A))$, a contradiction. Enlarge $b$ to cover the finitely many
roots of $P_v$ and obtain $\deg f(t^n)=O(n)$ for $n\ge0$.

Bell and Nguyen \cite[Theorem 3.1]{BN} prove that a function satisfying
\eqref{eq:cp} is a polynomial map if $\deg f(U^n)=O(n)$ for some
$U\in R$ with $U'\ne0$. Apply that theorem with $U=t$. Finally,
minimality of $\ell$ gives $q^\ell<16q$, and hence
$c_q>1/(256q^2)$.
\end{proof}

The gain comes from \eqref{eq:digitheight}: $q^{M-\ell}$ independent
polynomials have value degrees $O(q^M)$ on $V_{M+1}$. Ordinary monomials
of comparable $X$-degree have value degrees of order $Mq^M$.
The rational denominators of the digit polynomials do not enter the
value-height estimate, while the factor $\Pi_{s-1}$ ensures the exact
congruences needed for vanishing.

\section{The sharp linear case}
Every $\F_q$-linear $f:R\to R$ has a unique locally finite expansion
\begin{equation}\label{eq:linearseries}
 f(A)=\sum_{i\ge0}a_iE_i(A),\qquad a_i\in R.
\end{equation}
Indeed, recursively set
$a_n=f(t^n)-\sum_{i<n}a_iE_i(t^n)$. Lemma \ref{lem:carlitz} makes
the coefficients integral. The resulting map agrees with $f$ on the
basis $1,t,t^2,\ldots$, and $E_i(A)=0$ for $i>\deg A$.

\begin{proposition}\label{prop:linearcriterion}
The function \eqref{eq:linearseries} preserves prime congruences if and
only if $\Pi_i\mid a_i$ for every $i\ge0$.
\end{proposition}
\begin{proof}
Fix $P$ of degree $d$. For $i<d$, $H_i$ is a $P$-unit, so $E_i$ is a
polynomial over $R_{(P)}$. If $P\mid a_i$ for every $i\ge d$, all higher
terms vanish modulo $P$, proving sufficiency.

Conversely, subtract $\sum_{i<d}a_iE_i$ from $f$. The result modulo $P$
is a linear function on $R/(P)$ vanishing at the basis
$1,t,\ldots,t^{d-1}$, and hence vanishing everywhere.
Evaluating successively at $t^d,t^{d+1},\ldots$ and using
$E_i(t^i)=1$ proves $P\mid a_i$ for every $i\ge d$.
Apply this at all primes.
\end{proof}

\begin{lemma}\label{lem:transform}
Write $\alpha_n=\deg a_n$, $\beta_n=\deg f(t^n)$, and
$S_n=(q^n-1)/(q-1)$. Then
\begin{equation}\label{eq:transform}
 \alpha_n\le S_n+\max_{0\le j\le n}(\beta_j-S_j).
\end{equation}
Consequently,
\begin{equation}\label{eq:limsup}
 \limsup_n\frac{\alpha_n}{q^n}
 \le\max\left\{\frac1{q-1},\limsup_n\frac{\beta_n}{q^n}\right\}.
\end{equation}
\end{lemma}
\begin{proof}
The recursion and \eqref{eq:Edegree} give
$\alpha_n\le\max\{\beta_n,\max_{i<n}(\alpha_i+(n-i)q^i)\}$.
Since $(n-i)q^i\le S_n-S_i$, induction proves \eqref{eq:transform}.
For the limiting assertion it suffices to consider a finite right side.
Choose $c$ strictly larger than it, and then choose
$c_0<c$ with $c_0>1/(q-1)$ such that eventually $\beta_j\le c_0q^j$.
The finitely many earlier $\beta_j-S_j$ have an upper bound $B$.
Equation \eqref{eq:transform} gives
\[
 \alpha_n\le\max\{S_n+B,c_0q^n\}
\]
for all large $n$. Divide by $q^n$ and let $c$ decrease to the asserted
bound. The identically zero map causes no exception with the convention
$\deg0=-\infty$.
\end{proof}

\begin{proof}[Proof of Theorem \ref{thm:additive}]
If $\limsup\beta_n/q^n<q/(q-1)$, Lemmas \ref{lem:primecount} and
\ref{lem:transform} give $\deg a_n<d_n$ for all large $n$.
Proposition \ref{prop:linearcriterion} forces $a_n=0$ there. Thus
\eqref{eq:linearseries} is a finite polynomial.

For sharpness define
\begin{equation}\label{eq:boundary}
 f_*(A)=\sum_{i\ge0}\Pi_iE_i(A).
\end{equation}
This map is well-defined, $\F_q$-linear, and preserves prime congruences.
Put $C=q/(q-1)$. For $\deg A=n$,
\[
 \deg f_*(A)\le\max_{0\le i\le n}\{d_i+(n-i)q^i\}\le Cq^n,
\]
because $d_i<Cq^i$ and $C+k\le Cq^k$ for $k\ge0$.
For $i<n$ the same bound is at most $(C+1)q^{n-1}$: the sequence
$(C+k)/q^k$ decreases for $k\ge1$.
At $A=t^n$ the term with $i=n$ has degree $d_n$, and
$d_n/q^n\to C>(C+1)/q$. Eventually it strictly dominates all other
terms. Therefore $\deg f_*(t^n)=d_n$ for all sufficiently large $n$.
This gives the asserted sharp limit and proves nonpolynomiality, since
a fixed polynomial map has value degrees $O(n)$ at $t^n$.
\end{proof}

Wagner's characterization \cite[Theorem 3.2]{W74} for the stronger
all-moduli condition uses
\[
 L_i=\prod_{r=1}^i(t^{q^r}-t)
\]
in place of $\Pi_i$. Proposition \ref{prop:linearcriterion} is the squarefree
counterpart. The boundary construction above and the degree estimate
address the prime-congruence growth condition of \cite{BN}.

\section{Newton coefficients and the remaining threshold}
Fix an ordering $\alpha_0=0,\alpha_1=1,\ldots,\alpha_{q-1}$ of $\F_q$.
For the base-$q$ expansion $j=\sum_i j_iq^i$, put
$A_j=\sum_i\alpha_{j_i}t^i$. The first $q^m$ entries form $V_m$ and
$A_{q^m}=t^m$. Define
\begin{equation}\label{eq:newton}
 B_0(X)=1,\qquad
 B_n(X)=\frac{\prod_{j<n}(X-A_j)}{\prod_{j<n}(A_n-A_j)}\quad(n\ge1).
\end{equation}

\begin{proposition}\label{prop:newton}
Every function $f:R\to R$ has a unique locally finite expansion
\[
 f(A)=\sum_{n\ge0}b_nB_n(A),\qquad b_n\in R.
\]
It preserves prime congruences if and only if
\begin{equation}\label{eq:newtoncriterion}
 \Pi_{\lfloor\log_qn\rfloor}\mid b_n\qquad(n\ge1).
\end{equation}
There is no restriction on $b_0$.
\end{proposition}
\begin{proof}
For $P$ of degree $d$, the denominator in \eqref{eq:newton} has valuation
\[
 \sum_{k\ge1}\left\lfloor\frac n{q^{kd}}\right\rfloor.
\]
Indeed, each aligned block of $q^{kd}$ consecutive entries runs through
every residue modulo $P^k$. In the final incomplete block preceding
$A_n$, none has its residue: their differences from $A_n$ are nonzero
and have degree less than $kd$. The numerator at any input has at least
the same valuation. Thus all $B_n$ are $R$-valued. Since
$B_n(A_n)=1$ and $B_n(A_i)=0$ for $i<n$, integral triangular recursion
gives the expansion and uniqueness.

Fix $P$ of degree $d$. Terms with $n<q^d$ have $P$-unit denominators.
If $P\mid b_n$ for $n\ge q^d$, the other terms vanish pointwise modulo
$P$, proving sufficiency. Conversely, subtract the initial sum with
$n<q^d$. The resulting congruence-preserving function vanishes on the
complete residue system $A_0,\ldots,A_{q^d-1}$, hence everywhere modulo
$P$. Successive evaluation at the remaining $A_n$ gives $P\mid b_n$
for $n\ge q^d$. Taking all primes proves \eqref{eq:newtoncriterion}.
\end{proof}

This statement concerns first powers of primes. It makes no assertion
that the same ordering characterizes preservation of all prime powers;
compare the local constructions in \cite[Theorem 1.11]{LS}.

The sharp general question is consequently a truncation problem: does
\eqref{eq:newtoncriterion}, together with the strict subcritical bound
$(1-\varepsilon)d_n$ on degree-$n$ inputs, force $b_n=0$ eventually?
A finite expansion is equivalent to a polynomial map. The next example
shows why value growth by itself cannot supply the needed coefficient
margin.

\begin{proposition}\label{prop:obstruction}
Let $\delta(A)$ be $1$ at $A=0$ and $0$ elsewhere. Its coefficients in
\eqref{eq:newton} satisfy, for $m\ge1$,
\[
 \deg b_{q^m}=\frac{q^{m+1}-q}{q-1}-m,
 \qquad \frac{\deg b_{q^m}}{d_m}\longrightarrow1.
\]
\end{proposition}
\begin{proof}
The leading coefficient of interpolation at $A_0,\ldots,A_n$ gives
\[
 b_n=\frac{\prod_{j<n}(A_n-A_j)}{\prod_{j=1}^n(-A_j)}.
\]
For $n=q^m$, the numerator has degree $mq^m$. The denominator has
degree $m+\sum_{r=0}^{m-1}r(q-1)q^r$. Subtraction gives the formula;
Lemma \ref{lem:primecount} gives the limit.
\end{proof}

The function $\delta$ fails \eqref{eq:cp}, since $\delta(0)=1$ and
$\delta(P)=0$. It is an obstruction to an estimate using growth alone,
not a counterexample to the sharp question. A proof at every constant
below $q/(q-1)$ must use the congruences and the general coefficient
transform together, or supply a different rigidity argument.

\section{Verification and further directions}
The proofs hold for every prime power and do not depend on finite
computations. The accompanying exact-arithmetic checker tests both
prime and extension fields. The digit-basis audit covers 8,570 literal
values over $\F_2,\F_3,\F_4,\F_9$, including integrality, degrees, and
congruences. It also checks 56,320 instances of the height and
propagation inequalities and 2,816 dimension inequalities. A separate
audit covers the additive criterion, its boundary example, and the
general Newton criterion with positive and negative examples.

The main remaining problem is the sharp arbitrary-function threshold.
The integer-valued auxiliary family suggests optimizing its value-height
profile and replacing the rectangular set of indices in
\eqref{eq:auxiliary} by a weighted set. Those changes do not by themselves
prove the sharp constant. Theorems \ref{thm:general} and
\ref{thm:additive}, and their coefficient formulations, constitute one
growth-rigidity result rather than separate solutions to the full question.

\begin{thebibliography}{9}
\bibitem{BN}
J. P. Bell and K. D. Nguyen,
\emph{An analogue of Ruzsa's conjecture for polynomials over finite fields},
J. Combin. Theory Ser. A \textbf{178} (2021), 105337.
\href{https://doi.org/10.1016/j.jcta.2020.105337}{doi:10.1016/j.jcta.2020.105337}.
Theorem and question numbering here refers to
\href{https://arxiv.org/abs/1910.08255v1}{arXiv:1910.08255v1} (2019).
\bibitem{LS}
X. Li and M. Sha,
\emph{Congruence preserving functions in the residue class rings of
polynomials over finite fields},
Finite Fields Appl. \textbf{61} (2020), 101604.
\href{https://doi.org/10.1016/j.ffa.2019.101604}{doi:10.1016/j.ffa.2019.101604};
\href{https://arxiv.org/abs/1807.02379v2}{arXiv:1807.02379v2}.
\bibitem{W74}
C. G. Wagner, \emph{Linear pseudo-polynomials over $\operatorname{GF}[q,x]$},
Arch. Math. \textbf{25} (1974), 385--390.
\href{https://web.math.utk.edu/~cwagner/papers/pseudo.pdf}{Author's copy}.
\bibitem{W76}
C. G. Wagner, \emph{Polynomials over $\operatorname{GF}(q,x)$ with
integral-valued differences}, Arch. Math. \textbf{27} (1976), 495--501.
\href{https://web.math.utk.edu/~cwagner/papers/gfqx.pdf}{Author's copy}.
\end{thebibliography}
\end{document}
